\chapter{Գնահատում և ընդունման չափանիշներ}
\label{ch:assessment}

\section{Ընդունման չափանիշներ}
Նախագիծը ընդունվում է պաշտպանության, եթե.
\begin{itemize}
  \item MVP-ն (բաժին~\ref{ch:scope}) աշխատում է և ցուցադրվում է AUP ZU3-ի վրա;
  \item runtime կոնֆիգուրացիան Web հարթակից փոխում է FPGA-ի վարքը առանց նոր bitstream-ի;
  \item յուրաքանչյուր RTL մոդուլ ունի testbench, նախագիծը անցնում է timing-ը;
  \item կոդը և հրահանգները ռեպոզիտորիայում են, և մեկ այլ անձ կարող է կրկնել կառուցումը։
\end{itemize}

\section{Գնահատման աղյուսակ}
\begin{rubric}[Գնահատման չափանիշներ (յուրաքանչյուր ուսանողի համար՝ իր մասով).]
  \criterion{Տեխնիկական արդյունք}{35\%}{MVP-ն աշխատում է}{MVP + չափված բնութագրեր (throughput, latency)}{MVP + ընդլայնում, մաքուր ճարտարապետություն}
  \criterion{Վերիֆիկացիա և թեստավորում}{20\%}{Հիմնական մոդուլները թեստավորված են}{Ինքնաստուգվող թեստեր բոլոր մոդուլների համար}{Ավտոմատացված թեստեր, ծածկույթի վերլուծություն}
  \criterion{Վերլուծություն և ըմբռնում}{15\%}{Բացատրում է իր որոշումները}{Համեմատում է այլընտրանքները}{FPGA և ծրագրային մշակման խորքային համեմատություն}
  \criterion{Աշխատանքի գործընթաց}{10\%}{Հանգրվանները մեծ մասամբ պահպանված են}{Կանոնավոր commit-ներ և հաշվետվություններ}{Ինքնուրույնություն, ակտիվ համագործակցություն}
  \criterion{Փաստաթուղթ և պաշտպանություն}{20\%}{Ամբողջական դիպլոմ}{Հստակ կառուցվածք, նկարներ, չափումներ}{Հրապարակման մակարդակի շարադրանք և ցուցադրում}
\end{rubric}

\notebox{Աղյուսակը լրացնում է, բայց չի փոխարինում համալսարանի պաշտոնական գնահատման կանոնները։}
